% ============================================================================
% PREAMBLE: Common LaTeX packages and configuration for the Wiki
% ============================================================================
% This file defines all packages, global settings, and utility commands.
% Include this at the top of wiki_main.tex
% ============================================================================

\documentclass[11pt, a4paper, oneside]{book}

% ============================================================================
% CORE PACKAGES
% ============================================================================

\usepackage[utf8]{inputenc}
\usepackage[T1]{fontenc}
\usepackage[english]{babel}

% ============================================================================
% MATHEMATICS & SYMBOLS
% ============================================================================

\usepackage{amsmath}
\usepackage{amssymb}
\usepackage{amsfonts}
\usepackage{mathtools}
% Removed: physics (conflicts with custom commands)
\usepackage{tensor}            % Tensor notation
\usepackage{bm}                % Bold math symbols

% ============================================================================
% LAYOUT & FORMATTING
% ============================================================================

\usepackage[margin=1.2in]{geometry}
\usepackage{fancyhdr}
\usepackage{setspace}
\onehalfspacing                % 1.5 line spacing

% Page style
\pagestyle{fancy}
\fancyhf{}
\fancyhead[LE]{\leftmark}
\fancyhead[RO]{\rightmark}
\fancyfoot[C]{\thepage}
\renewcommand{\headrulewidth}{0.5pt}

% ============================================================================
% GRAPHICS & FIGURES
% ============================================================================

\usepackage{graphicx}
\usepackage{tikz}
\usepackage{pgfplots}
\pgfplotsset{compat=1.18}
\usepackage{subcaption}
\usepackage{float}

% ============================================================================
% TABLES
% ============================================================================

\usepackage{booktabs}
\usepackage{array}
\usepackage{multirow}
\usepackage{dcolumn}           % Align on decimal point

% ============================================================================
% CROSS-REFERENCING & HYPERLINKS
% ============================================================================

\usepackage{hyperref}
\hypersetup{
    colorlinks   = true,
    linkcolor    = blue,
    urlcolor     = blue,
    citecolor    = red,
    bookmarks    = true,
    bookmarksnumbered = true,
    pdfauthor    = {Tirocinio MPMSSIA},
    pdftitle     = {3D Power Spectrum \& Weak Lensing Wiki},
}

\usepackage{nameref}           % Use \nameref for cross-references
\usepackage[noabbrev]{cleveref}  % Smart reference formatting (\cref, \Cref)
\usepackage{bookmark}          % Better PDF bookmarks

% Setup cref names
\crefname{equation}{Eq.}{Eqs.}
\crefname{figure}{Fig.}{Figs.}
\crefname{table}{Tab.}{Tabs.}
\crefname{section}{Sec.}{Secs.}

% ============================================================================
% BIBLIOGRAPHY
% ============================================================================

\usepackage[
    style=numeric-comp,
    sorting=nyt,
    backend=bibtex,
]{biblatex}
\addbibresource{_common/bib.bib}

% ============================================================================
% CODE & LISTINGS
% ============================================================================

\usepackage{listings}
\usepackage{xcolor}

\lstset{
    basicstyle=\ttfamily\small,
    keywordstyle=\color{blue},
    commentstyle=\color{gray},
    stringstyle=\color{red},
    breaklines=true,
    postbreak=\mbox{\textcolor{gray}{$\hookrightarrow$}\space},
    frame=single,
    framerule=0.4pt,
    numberstyle=\tiny\color{gray},
    numbersep=8pt,
}

% ============================================================================
% THEOREM ENVIRONMENTS
% ============================================================================

\usepackage{amsthm}
\usepackage{thmtools}

\theoremstyle{definition}
\newtheorem{definition}{Definition}[chapter]
\newtheorem{example}{Example}[chapter]
\newtheorem{exercise}{Exercise}[chapter]

\theoremstyle{plain}
\newtheorem{theorem}{Theorem}[chapter]
\newtheorem{lemma}[theorem]{Lemma}
\newtheorem{proposition}[theorem]{Proposition}
\newtheorem{corollary}[theorem]{Corollary}

\theoremstyle{remark}
\newtheorem{remark}{Remark}[chapter]
\newtheorem{note}{Note}[chapter]

% ============================================================================
% MISCELLANEOUS
% ============================================================================

\usepackage{xspace}            % Smart spacing after macros
\usepackage{comment}           % Comment environments
\usepackage[final]{pdfpages}   % Include external PDFs if needed
\usepackage{chemfig}           % Chemical structures (if needed)

% ============================================================================
% TABLE OF CONTENTS SETTINGS
% ============================================================================

\setcounter{tocdepth}{2}       % Show chapters and sections in TOC
\setcounter{secnumdepth}{3}    % Number up to subsubsections

% ============================================================================
% CUSTOM TITLE PAGE
% ============================================================================

\title{\textbf{3D Power Spectrum \& Weak Lensing Theory}\\[0.5em]
       {\Large A Comprehensive Wiki with Spherical Fourier-Bessel Framework}}
\author{Tirocinio MPMSSIA Internship Project}
\date{\today}

% ============================================================================
% END PREAMBLE
% ============================================================================
